% arara: lualatex
% arara: lualatex
\documentclass{article}
% \animategraphics[width=5cm,loop,autoplay]{10}{path-to-file}{0}{118}
\PassOptionsToPackage{colorlinks = true}{hyperref}
\usepackage{
    pline-vec
    ,minted
    ,amssymb
    ,caption
    ,bookmark
    ,hyperref
}
\title{pline-vec}
\author{Jasper}
\date{\today}
\begin{document}
\maketitle
\tableofcontents
% See https://tex.stackexchange.com/a/31787/319072
\section{Introduction}
pline-vec is a 3D toolkit for Ti\textit{k}Z. This document will
one day be the documentation for pline-vec, though it currently
serves as a test of my functions.
\section{Examples}

\begin{tikzpicture}
    \pvvector[red,points=both]{0,0}{1,0}
    \pvvector{0,1}{1,1}
\end{tikzpicture}

\subsection{The pline-vec math libraries}
\subsubsection{Parametric mappings}
\subsection{Setting the main orientation}
\begin{tikzpicture}[main orientation = {pi/2}{pi/3}{5*pi/6}]
    \draw[->] (0,0) -- (1,0,0) node[pos = 1.2]{\(x\)};
    \draw[->] (0,0) -- (0,1,0) node[pos = 1.2]{\(y\)};
    \draw[->] (0,0) -- (0,0,1) node[pos = 1.2]{\(z\)};
\end{tikzpicture}
\subsubsection{Setting sub-orientations}
\begin{tikzpicture}[main orientation = {pi/2}{pi/3}{5*pi/6}]
    \draw[->] (0,0) -- (1,0,0) node[pos = 1.2]{\(x\)};
    \draw[->] (0,0) -- (0,1,0) node[pos = 1.2]{\(y\)};
    \draw[->] (0,0) -- (0,0,1) node[pos = 1.2]{\(z\)};

    \draw
        [->,red] 
        (0,0) 
        -- 
        (
            {ZYZx(-pi/6,pi/6,pi/6,1,0,0)}
            ,{ZYZy(-pi/6,pi/6,pi/6,1,0,0)}
            ,{ZYZz(-pi/6,pi/6,pi/6,1,0,0)}
        ) 
        node
        [pos = 1.2]
        {\(x\)}
        ;
    \draw
        [->,red] 
        (0,0) 
        -- 
        (
            {ZYZx(-pi/6,pi/6,pi/6,0,1,0)}
            ,{ZYZy(-pi/6,pi/6,pi/6,0,1,0)}
            ,{ZYZz(-pi/6,pi/6,pi/6,0,1,0)}
        ) 
        node
        [pos = 1.2]
        {\(y\)}
        ;
    \draw
        [->,red] 
        (0,0) 
        -- 
        (
            {ZYZx(-pi/6,pi/6,pi/6,0,0,1)}
            ,{ZYZy(-pi/6,pi/6,pi/6,0,0,1)}
            ,{ZYZz(-pi/6,pi/6,pi/6,0,0,1)}
        ) 
        node
        [pos = 1.2]
        {\(z\)}
        ;
\end{tikzpicture}
\subsection{Rendering parametric curves}
\begin{tikzpicture}[main orientation = {pi/2}{pi/3}{5*pi/6}]
    \pvappendcurve[
        u min = 0
        ,u max = 2*pi
        ,u samples = 200
    ]{3*math.cos(u) + 1.5*pv_sphere_x(u,3*u)}
    {3*math.sin(u) + 1.5*pv_sphere_y(u,3*u)}
    {1.5*pv_sphere_z(u,3*u)}
    \pvappendsurface[
        u min = pi/2
        ,u max = 2*pi
        ,u samples = 36
        ,v min = 0
        ,v max = 2*pi
        ,v samples = 36
    ]{3*math.cos(u) + pv_sphere_x(u,v)}
    {3*math.sin(u) + pv_sphere_y(u,v)}
    {pv_sphere_z(u,v)}
    \pvappendsurface[
        u min = 0
        ,u max = 2*pi
        ,u samples = 36
        ,v min = 0
        ,v max = 2*pi
        ,v samples = 9
    ]{3*math.cos(u) + 0.3*pv_sphere_x(u,v)}
    {3*math.sin(u) + 0.3*pv_sphere_y(u,v)}
    {0.3*pv_sphere_z(u,v)}
    \pvappendsolid[
    u min = pi/2,
    u max = 2*pi,
    u samples = 30,
    v min = -pi/2,
    v max = pi/2,
    v samples = 30,
    w min = 1,
    w max = 1.2,
    w samples = 3
    ]{w*pv_sphere_x(u,v)}
    {w*pv_sphere_y(u,v)}
    {w*pv_sphere_z(u,v)+3}
    \pvrendersegments
    \draw[->] (0,0) -- (1,0,0) node[pos = 1.2]{\(x\)};
    \draw[->] (0,0) -- (0,1,0) node[pos = 1.2]{\(y\)};
    \draw[->] (0,0) -- (0,0,1) node[pos = 1.2]{\(z\)};
\end{tikzpicture}
\subsection{Rendering parametric surfaces}
\subsection{Rendering parametric solids}
\subsection{
    Rendering parametric curves, 
    surfaces and solids all at once
}
\subsubsection{The problem with tangency}
\subsection{Clipping planes by a rectangular prism}
\begin{tikzpicture}[main orientation = {pi/2}{pi/3}{5*pi/6}]
    \pvgetrectangularprism
    \backfaces
    \pvgetplane{thename}
    \draw[spath/use = planethename,preaction = {fill = gray}];
    \frontfaces
\end{tikzpicture}
\subsection{Intersecting planes --- in progress}
\end{document}